\documentclass[11pt,a4paper]{article}
\usepackage{curso}
\begin{document}
\portada{Lección}{5.1}{Por qué el hardware digital es ineficiente para la IA probabilística}

\noindent\textbf{Unidad 5.}\quad\textbf{Objetivos:} desglosar la energía de una actualización de Gibbs en digital; entender por qué el muestreo está limitado
por memoria y por qué los márgenes de ruido cuestan; estimar la brecha frente a un pbit físico.
\textbf{Prerrequisitos:} 0.1 (energía y roofline), 4.3 y 4.5 (Gibbs y pbits).

\section{Anatomía de una actualización}
Una actualización de un espín de grado 16 lee 16 estados y 16 acoplamientos, suma el campo local, evalúa una sigmoide, genera un número aleatorio,
compara y escribe. La aritmética es mínima ($\sim30$ operaciones); el costo está en mover datos. Con la tabla de 45 nm (lección 0.1): suma FP
$\approx0{,}9$ pJ, lectura SRAM de 32 bits $\approx5$ pJ, lectura DRAM $\approx640$ pJ.

\section{Muestreo limitado por memoria}
Con $\sim1$ FLOP por byte, la actualización está muy a la izquierda del codo del roofline: las unidades aritméticas esperan datos. El grafo disperso
de un modelo de Ising tampoco permite la reutilización en caché de una multiplicación densa.

\section{El costo de la aleatoriedad y de los márgenes}
Los generadores pseudoaleatorios (Philox, xorshift) gastan varias operaciones por palabra. Un generador de solo transistores basado en ruido de disparo
produce un bit aleatorio con cientos de attojulios (pensum, Unidad 5). Además, cada bit digital se guarda con barreras de decenas de $\kB T$
(lección 0.2): se paga por eliminar el ruido y luego por fabricarlo.

\begin{ejemplo}[ (estimación de orden de magnitud \evE)]
Digital: $12\times5$ pJ (SRAM) $+\ 3$ pJ (PRNG) $+\ 5$ pJ (ALU) $\approx70$ pJ por actualización. Z1: $1\ \text{W}/(269\,568\times50\ \text{MHz})\approx74$ fJ.
Cociente $\approx10^3$, sin contar la memoria externa. Ambas cifras son cotas o estimaciones; la comparación justa es por muestra independiente y a nivel
de sistema (lección 5.8).
\end{ejemplo}

\begin{aplicacion}
Diseño: acoplamientos junto a cada pbit (cómputo en memoria), ruido físico en lugar de PRNG, lecturas poco frecuentes. Laboratorio L7: medir J/muestra de
Gibbs en GPU y compararlo con la estimación para pbits.
\end{aplicacion}

\begin{error}
\begin{itemize}[nosep]
\item Comparar FLOPs en lugar de bytes movidos. \quad Comparar energía de tejido de un chip con energía de sistema de una GPU.
\item Ignorar costo y calidad del PRNG. \quad Tratar estimaciones \evE{} como mediciones.
\end{itemize}
\end{error}

\begin{thebibliography}{9}\small
\bibitem[21]{e1} Extropic, Thermodynamic Computing From Zero to One, 2025.
\bibitem[22]{e2} Extropic, From One to One Billion: Torx, Thermalizers, and Z1, 2026.
\bibitem{horowitz} M. Horowitz, Computing's energy problem (and what we can do about it), ISSCC 2014.
\end{thebibliography}
\end{document}
